\documentclass[12pt, a4paper]{article}

% --- PACOTES ESSENCIAIS ---
\usepackage[utf8]{inputenc}      % Permite acentuação direta
\usepackage[brazilian]{babel}    % Traduz e ajusta hifenação para Português-BR
\usepackage[T1]{fontenc}         % Melhora a fonte e a hifenação
\usepackage{microtype}           % Para micro-ajustes tipográficos e justificação perfeita
\usepackage{ragged2e}            % Para alinhamento robusto (inclui \justifying)
\usepackage{listings}            % Para snippets de código


% --- LAYOUT E DESIGN ---
\usepackage[margin=2.5cm]{geometry} % Define as margens (2.5cm)
\usepackage{graphicx}            % Para incluir imagens (ex: logo)
\usepackage{fancyhdr}            % Para cabeçalhos e rodapés
\usepackage[svgnames]{xcolor}    % Para cores
\usepackage{titlesec}            % Para customizar seções


% --- TABELAS ---
\usepackage{booktabs}            % Tabelas mais bonitas (toprule, midrule, bottomrule)
\usepackage{longtable}           % Para tabelas que quebram por várias páginas
\usepackage{array}               % Para customizar colunas de tabelas
\usepackage[table]{xcolor}       % Para cores em tabelas (ex: \rowcolor)
\usepackage{tabularx}            % Para tabelas com largura definida

% --- DESIGN --- 
\usepackage{tikz}                 % Pacote para desenhar
\usetikzlibrary{calc}
\definecolor{AzulCapa}{RGB}{0, 61, 110} % Um tom de azul escuro e profissional
\definecolor{AzulCapa2}{HTML}{1A4E8A} % <-- Note: modelo HTML e SEM '#'
% --- LISTAS ---
\usepackage{enumitem}            % Para customizar listas

% --- CONFIGURAÇÃO DE CABEÇALHO E RODAPÉ (Exemplo Elegante) ---
\pagestyle{fancy}
\fancyhf{} % Limpa todos os campos
\lhead{\small Laboratório de Ensaio acreditado pela CGCRE de acordo com a ABNT NBR ISO/IEC 17025, sob o número CRL 1227}
\cfoot{\small Página \thepage}
\renewcommand{\headrulewidth}{0.4pt}
\renewcommand{\footrulewidth}{0.4pt}

% --- CONFIGURAÇÃO DO PACOTE LISTINGS (para snippets de código) ---
\definecolor{codegray}{rgb}{0.5,0.5,0.5} % Cor cinza para números
\definecolor{backcolor}{rgb}{0.98,0.98,0.98} % Cor de fundo leve

\lstdefinestyle{mystyle}{
    backgroundcolor=\color{backcolor},   
    basicstyle=\ttfamily\footnotesize, % Fonte (Typewriter) e tamanho
    breakatwhitespace=false,         
    breaklines=true,                 % Quebra linhas longas
    captionpos=b,                    % Posição da legenda (bottom)
    frame=single,                    % Borda simples ao redor
    keepspaces=true,                 
    numbers=left,                    % Números de linha à esquerda
    numbersep=5pt,                   % Espaço entre número e código
    numberstyle=\tiny\color{codegray}, % Estilo dos números de linha
    showspaces=false,                
    showstringspaces=false,
    showtabs=false,                  
    tabsize=2
}

\lstset{style=mystyle} % Aplica este estilo como padrão
% --- FIM DA CONFIGURAÇÃO DE LISTINGS ---

% --- CUSTOMIZAÇÃO DE TÍTULOS (Opcional) ---
\titleformat{\section}
  {\normalfont\Large\bfseries\color{SteelBlue}} % Fonte e cor das seções
  {\thesection}
  {1em}
  {}

\titleformat{\subsection}
  {\normalfont\large\bfseries} % Fonte das subseções
  {\thesubsection}
  {1em}
  {}

% --- INFORMAÇÕES DO DOCUMENTO ---
\title{Relatório de Ensaio de Módulo Criptográfico \vspace{0.5cm} \\ 
       \large \textit{Baseado no modelo LASPI-F-033 v.003}}
       
\author{Laboratório de Aplicações Tecnológicas (LASPI - UFRJ)}
\date{XX de Mês de ANO} % Ou use \today para a data atual


% %%%%%%%%%%%%%%%%%%%%%%%%%%%%%%%%%%%%%%%%%%%%%%%%%%%%%%%
% --- INÍCIO DO DOCUMENTO ---
% %%%%%%%%%%%%%%%%%%%%%%%%%%%%%%%%%%%%%%%%%%%%%%%%%%%%%%%
\begin{document}

% %%%%%%%%%%%%%%%%%%%%%%%%%%%%%%%%%%%%%%%%%%%%%%%%%%%%%%%
% --- PÁGINA DE CAPA (Estilo Diagramação) ---
% %%%%%%%%%%%%%%%%%%%%%%%%%%%%%%%%%%%%%%%%%%%%%%%%%%%%%%%
% %%%%%%%%%%%%%%%%%%%%%%%%%%%%%%%%%%%%%%%%%%%%%%%%%%%%%%%
% --- PÁGINA DE CAPA (Design com Diagonal e Semicírculo) ---
% %%%%%%%%%%%%%%%%%%%%%%%%%%%%%%%%%%%%%%%%%%%%%%%%%%%%%%%
\begin{titlepage}
    \thispagestyle{empty} % Remove cabeçalho/rodapé

    \begin{tikzpicture}[remember picture, overlay]

    % --- CAMADA 1: Fundo (AzulCapa - O mais escuro) ---
        \fill[AzulCapa] (current page.south west) rectangle (current page.north east);

        % --- CAMADA 2: Divisão Diagonal (AzulSwoosh - O mais claro) ---
        % (Desenha um polígono da borda direita para a esquerda)
        \fill[AzulCapa2] 
            ([xshift=-10cm] current page.north east) -- % Ponto 1 (Topo-Direita)
            (current page.north east) --                % Ponto 2 (Canto Sup-Dir)
            (current page.south east) --                % Ponto 3 (Canto Inf-Dir)
            ([yshift=20cm] current page.south west) --  % Ponto 4 (Baixo-Esquerda)
            cycle; % Fecha o polígono

    % --- CAMADA 3: Faixa Branca com Semicírculo ---
        % (Define as dimensões)
        \def\StripeRadius{1.75cm}  % Raio do semicírculo (Metade da largura)
        \def\StripeXCenter{18.0cm} % Posição X da faixa (2.5cm + 1.75cm)
        \def\StripeYBottom{26.5cm} % Posição Y do centro do arco (de baixo para cima)
        % --- !!! AJUSTE AQUI A ALTURA DA FAIXA !!! ---
        % (Altura medida a partir da base da página)
        \def\StripeYTop{26.0cm} % <-- Mude este valor para encurtar
        
        % (Coordenada do centro do arco)
        \coordinate (ArcCenter) at ([xshift=\StripeXCenter, yshift=\StripeYBottom] current page.south west);

        % (Desenha a faixa em DUAS PARTES - Correção da Altura)
        
        % 1. O Retângulo (com altura definida)
        \fill[AzulCapa2] 
        % Canto Inferior-Esquerdo do retângulo
        ([xshift=\StripeXCenter-\StripeRadius, yshift=\StripeYBottom] current page.north west)
            rectangle % Desenha até o...
        % Canto Superior-Direito do retângulo
        ([xshift=\StripeXCenter+\StripeRadius, yshift=\StripeYTop] current page.south west);

        % 2. O Semicírculo (preenchido)
        \fill[AzulCapa2] (ArcCenter) -- ++(0:\StripeRadius) arc (0:-180:\StripeRadius) -- cycle;
        
     % --- CAMADA 4: Logo (Sobre o semicírculo) ---
        % (Coloca o logo no centro do arco)
        \node at ([yshift=1.25cm]ArcCenter) {
         % Use seu logo com fundo transparente aqui!
         % O logo-2.png (preto, transparente) que você enviou funciona bem.
         \includegraphics[width=1.5cm]{assets/logo-teste.png}
         };

    % --- CAMADA 5: TEXTO (Lado Direito) ---
        % (Exatamente o mesmo código de texto de antes)

        % --- Título ---
        \node[text=white,
            anchor=north west,
            text width=11.5cm,
            align=left]
            at ([xshift=7cm, yshift=-8cm] current page.north west)
        {
            \normalfont\huge\bfseries
            Relatório de Ensaios Técnicos Módulo Criptográfico padrão ICP-Brasil
        
            \vspace{1.5cm}
            
            \large\normalfont\itshape

        };

        % --- Autor (LASPI) ---
        \node[text=white,
            anchor=north west,
            text width=11.5cm,
            align=left]
            at ([xshift=7cm, yshift=-13cm] current page.north west)
        {
        \Large
        Laboratório de Aplicações Tecnológicas  \newline para  o Setor Produtivo Industrial 
        };

        % --- UFRJ ---
        \node[text=white,
            anchor=south west,
            text width=11.5cm,
            align=left]
            at ([xshift=5.5cm, yshift=3cm] current page.south west)
        {
        \Large
         Universidade Federal do Rio de Janeiro
        };

    \end{tikzpicture}
\end{titlepage}

% %%%%%%%%%%%%%%%%%%%%%%%%%%%%%%%%%%%%%%%%%%%%%%%%%%%%%%%
% --- FIM DA PÁGINA DE CAPA ---
% %%%%%%%%%%%%%%%%%%%%%%%%%%%%%%%%%%%%%%%%%%%%%%%%%%%%%%%


% --- Prepara o resto do documento ---
\cleardoublepage % Garante que o conteúdo comece na próxima página ímpar
\setcounter{page}{1} % Reinicia a contagem de página
\pagestyle{fancy}   % Retoma o estilo de página normal (cabeçalho/rodapé)


% %%%%%%%%%%%%%%%%%%%%%%%%%%%%%%%%%%%%%%%%%%%%%%%%%%%%%%%
% --- SEGUNDA PÁGINA (FOLHA DE ROSTO) ---
% %%%%%%%%%%%%%%%%%%%%%%%%%%%%%%%%%%%%%%%%%%%%%%%%%%%%%%%
\begin{titlepage}
    \thispagestyle{empty} % Sem cabeçalho/rodapé nesta página

    
\vfill % Adiciona espaço vertical flexível ANTES do bloco

    % %%%%%%%%%%%%%%%%%%%%%%%%%%%%%%%%%%%%%%%%%%%%%%%%%%%%%%%
    % --- BLOCO TÍTULO E LOGO (LADO-A-LADO) ---
    % %%%%%%%%%%%%%%%%%%%%%%%%%%%%%%%%%%%%%%%%%%%%%%%%%%%%%%%
    
    % Centraliza o "bloco" inteiro (Título + Logo) na página
    \noindent\hspace{\stretch{1}}%
    
    % --- TÍTULO ---
    % ( \begin{center} está correto aqui, pois o texto deve ser centralizado)
    \begin{center}
    {\Large\bfseries Relatório de Ensaios Técnicos \par}
     \vspace{0.5cm} % Pequeno espaço
     {\large\bfseries Módulo Criptográfico padrão ICP-Brasil \par}
            
            % --- PONTO DE ÂNCORA PARA O LOGO ---
            \tikz[remember picture, overlay] \coordinate (centroTitulo);
    \end{center}
    
    % %%%%%%%%%%%%%%%%%%%%%%%%%%%%%%%%%%%%%%%%%%%%%%%%%%%%%%%
    % --- LOGO FLUTUANTE (TIKZ) ---
    % %%%%%%%%%%%%%%%%%%%%%%%%%%%%%%%%%%%%%%%%%%%%%%%%%%%%%%%
    \begin{tikzpicture}[remember picture, overlay]
        % 'overlay' não ocupa espaço
        % 'remember picture' usa a âncora que definimos
        
        \node[anchor=west] % Alinha o logo pela borda esquerda (west)
              at ([xshift=6cm, yshift=1cm]centroTitulo) % Posição: 6cm à direita da âncora
        {
            \includegraphics[width=0.15\textwidth]{assets/crl.png}
        };
    \end{tikzpicture}

    \vfill % Adiciona espaço vertical flexível DEPOIS do bloco
    
    % --- PARÁGRAFO DE DESCRIÇÃO ---
    % --- SOLUÇÃO NOVA AQUI ---
    \noindent\hspace{\stretch{1}}%
    \begin{parbox}{0.9\textwidth}
        \small % Fonte pequena para o parágrafo
        % (Sem \centering, sem \justifying. Apenas o padrão.)
        Relatório de ensaios técnicos para verificação de requisitos em Módulo Criptográfico submetido a processo de certificação, em conformidade com as Portarias INMETRO nº 130, de 19 de março de 2021, INMETRO nº 200, de 29 de abril de 2021 e Manuais de Condutas Técnicas (MCT), do Instituto Nacional de Tecnologia da Informação – ITI.
    \end{parbox}%
    \hspace{\stretch{1}}
    
    \vfill % Adiciona espaço vertical flexível
    
    % --- NÚMERO SEQUENCIAL ---
    % ( \begin{center} está correto aqui)
    \begin{center} 
        \textbf{Número sequencial do Relatório de Ensaios Técnicos:} X
    \end{center}
    
    \vfill % Adiciona espaço vertical flexível
    
    % --- TABELA DE INFORMAÇÕES ---
    % (Correto, não mexe. Já ocupa a largura total)
    \begin{tabularx}{\textwidth}{ >{\bfseries}r X }
        Órgão Emitente: & 
        ``Laboratório de Aplicações Tecnológicas para o Setor Produtivo e Industrial - LASPI, da Escola Politécnica, POLI/UFRJ, CNPJ 33.663.683/0006-20, Av. Athos da Silveira Ramos, Nº149, Centro de Tecnologia, Bloco H, Sala 108, Ilha da Cidade Universitária, Campus Fundão, Rio de Janeiro, RJ, CEP 21941-909.'' \\
        \addlinespace[10pt]
        Organismo de Certificação: & 
        ``NCC CERTIFICAÇÕES DO BRASIL LTDA, Sociedade Empresária Limitada, com sede na Av. Orosimbo Maia, 360 Sala 112. Centro, Campinas, SP, Brasil, CEP: 13.010-211.'' \\
        \addlinespace[10pt]
        Fabricante: & XXXX \\
        \addlinespace[10pt]
        Número do processo no OCP: & ... \\ 
        \addlinespace[10pt]
        Data de recebimento das amostras: & ... \\
        \addlinespace[10pt]
        Modelo (família) do módulo: & ... \\
    \end{tabularx}
    
    \vfill % Adiciona espaço vertical flexível
    
    % --- PARÁGRAFO DE DESCRIÇÃO 2 ---
    % --- SOLUÇÃO NOVA AQUI ---
    \noindent\hspace{\stretch{1}}%
    \begin{parbox}{0.9\textwidth}
        \small 
        Os resultados descritos neste relatório referem-se apenas à(s) amostra(s) ensaiada(s). 
        A fim de atestar a veracidade das informações contidas neste documento e torná-lo válido, o representante do órgão emitente firma o presente Relatório de Ensaios Técnicos.
    \end{parbox}%
    \hspace{\stretch{1}}
        
    \vfill % Adiciona espaço vertical flexível

    % %%%%%%%%%%%%%%%%%%%%%%%%%%%%%%%%%%%%%%%%%%%%%%%%%%%%%%%
    % --- BLOCO DE ASSINATURA ---
    % %%%%%%%%%%%%%%%%%%%%%%%%%%%%%%%%%%%%%%%%%%%%%%%%%%%%%%%
    \vfill % Adiciona um espaço flexível *antes* da assinatura
    
    % ( \begin{center} está correto aqui)
    \begin{center} 
        \begin{minipage}{0.8\textwidth}
            \centering % (Este \centering interno é para o conteúdo da assinatura)
            \vspace{1.5cm} 
            \rule{10cm}{0.4pt} \\
            \textbf{Carlos José Ribas d´Ávila} \\
            \textit{Responsável Técnico / Coordenador do LASPI} \\
            \vspace{1cm}
            \textbf{Rio de Janeiro, XX de XX de XX}
        \end{minipage}
    \end{center}
    
    \vfill % Espaço final para balancear

\end{titlepage}
% %%%%%%%%%%%%%%%%%%%%%%%%%%%%%%%%%%%%%%%%%%%%%%%%%%%%%%%
% --- FIM DA SEGUNDA PÁGINA ---
% %%%%%%%%%%%%%%%%%%%%%%%%%%%%%%%%%%%%%%%%%%%%%%%%%%%%%%%

% %%%%%%%%%%%%%%%%%%%%%%%%%%%%%%%%%%%%%%%%%%%%%%%%%%%%%%%
% --- SUMÁRIO E LISTAS ---
% %%%%%%%%%%%%%%%%%%%%%%%%%%%%%%%%%%%%%%%%%%%%%%%%%%%%%%%
\tableofcontents
\cleardoublepage % Força o início na próxima página (ímpar)

\listoffigures   
\cleardoublepage % Força o início na próxima página (ímpar)

\listoftables    
\cleardoublepage % Força o início na próxima página (ímpar)

\lstlistoflistings 
\cleardoublepage % Força o início na próxima página (ímpar)

\newpage         % Garante que o conteúdo principal comece na próxima página

% --- SEÇÃO DE APRESENTAÇÃO ---
\section*{Apresentação}
\addcontentsline{toc}{section}{Apresentação} % Adiciona ao Sumário (se houver)

 O presente documento descreve os resultados do processo de avaliação do módulo criptográfico XXX, da fabricante XXX, para verificação de atendimento aos requisitos de conformidade presentes no Manual de Condutas Técnicas 7, Volume II, do Instituto Nacional de Tecnologia da Informação – ITI para o nível NSH X. 

% --- SEÇÃO DO PROCESSO DE AVALIAÇÃO ---
\section*{O processo de avaliação}
\addcontentsline{toc}{section}{O processo de avaliação}

 A avaliação foi realizada no período de XX de MÊS a XX de MÊS de ANO, nas dependências do Laboratório de Aplicações Tecnológicas para o Setor Produtivo Industrial da Universidade Federal do Rio de Janeiro - LASPI.

A descrição detalhada dos testes é apresentada a seguir.

% --- TABELAS INICIAIS ---
\subsection*{Versões Utilizadas na Avaliação}

% 'h!' tenta forçar a tabela a ficar "aqui"
\begin{table}[h!]
\centering
\begin{tabular}{ll}
\toprule
\textbf{Componente} & \textbf{Descrição} \\
\midrule
 Fabricante & XX\\
 Modelo     & XX \\  
 Hardware   & XX \\  
 Software   & XX \\  
 Firmware   & XX \\  
\bottomrule
\end{tabular}
\caption{Versões dos componentes avaliados.}
\end{table}



% --- TABELA DE VERSÃO E CONDIÇÕES AMBIENTAIS ---
\begin{table}[h!]
\centering
\begin{tabular}{ll}
\toprule
\multicolumn{2}{l}{\textbf{Histórico de Versões do Relatório}} \\
\midrule
Versão 1.0 & Versão revisada pelo LASPI. \\
% Adicione mais versões se necessário
\bottomrule
\addlinespace[10pt] % Adiciona um espaço vertical
\toprule
\multicolumn{2}{l}{\textbf{Condições Ambientais dos Ensaios}} \\
\midrule
Temperatura & 21 ºC - 23 ºC \\
Umidade     & 47\% - 48\% \\
\bottomrule
\end{tabular}
\caption{Metadados do Relatório e Ambiente.}
\end{table}

\newpage

% --- SEÇÃO PRINCIPAL DE REQUISITOS (NOVO DESIGN) ---
\section{Resultados dos Ensaios}

% --- CABEÇALHO ÚNICO (NÃO SE REPETE MAIS) ---
% Centraliza o título mestre da seção
\begin{center}
    {\large\textbf{Requisitos de Especificação do Módulo Criptográfico conforme}} \\
    {\small LASPI-PR-008 Procedimento de Ensaios Técnicos para ICP-BRASIL v.010}
\end{center}
\vspace{0.5cm} % Adiciona um espaço
% --- FIM DO CABEÇALHO ÚNICO ---




\subsection{REQUISITO III.1.1} 
    
    % --- Tabela minimalista (Design Limpo) ---
    
    \begin{tabularx}{\textwidth}{l X} 
    \toprule
    
    \textbf{Referência Normativa} & \textbf{Requisito e Item} \\ 
    \midrule
    
    \rowcolor{gray!15} % Cor de fundo
    
    \textbf{MCT 7 Vol. II} & 
    "A parte interessada deve fornecer documentação específica dos componentes de hardware, software e firmware do módulo criptográfico além da fronteira criptográfica que delimita tais componentes." \\
    \bottomrule
    \end{tabularx}
    % --- Fim da Tabela ---
    
    \vspace{1cm} 
    
    % --- Procedimentos (Texto Livre) ---
    \subsubsection*{Procedimentos de ensaio para NSH 1, 2 e 3:}
    \begin{itemize}[leftmargin=*, nosep]
        \item \textbf{EN.III.1.1.01:} Verificar se a documentação inclui uma “lista de componentes principais” que descreve todos componentes de hardware, software e firmware do módulo criptográfico. Verificar se a “lista de componentes principais” inclui, mas não se limita aos seguintes tipos de componentes:
        \begin{itemize}
            \item Processadores, incluindo microprocessadores, processadores de sinal digital, processadores  personalizados/dedicados, microcontroladores, ou quaisquer outros tipos de processadores;
            \item Circuitos integrados de memória ROM (Read Only Memory) para código executável de programas e dados [tal item pode incluir MPROM (Mask-Programmed ROM), PROM (Programmable ROM), EPROM (Erasable PROM), EEPROM (Electrically Erasable PROM) ou FLASH];
            \item Circuitos integrados de memória RAM (Random Access Memory) para armazenamento de dados temporários;
            \item Circuitos integrados semi dedicado (semi-custom) de aplicação específica, tais como, gate arrays, programmable logic arrays, field programmable gate arrays, ou outros dispositivos lógicos programáveis;
            \item Circuitos integrados totalmente dedicados (fully custom) e de aplicação específica, incluindo quaisquer circuitos integrados criptográficos e dedicados;
            \item Outros elementos ativos de circuito eletrônico (a documentação não deve listar elementos passivos de circuito eletrônico, tais como, resistores pull up/pull down ou capacitores bypass se eles não desempenharem um papel significativo na segurança do módulo criptográfico e não estiverem na fronteira criptográfica;
            \item Componentes de fornecimento de energia, incluindo alimentação (power supply), módulos de conversão de voltagem (por exemplo: módulos AC-DC ou DC-DC), transformadores, conectores de entrada de energia e conectores de saída de energia;
            \item Placa de circuito impresso ou outras superfícies de montagem de componentes;
            \item Encapsulamentos/revestimentos, incluindo quaisquer portas de acesso removíveis ou coberturas/camadas;
            \item Conectores físicos para dispositivos externos ao módulo criptográfico, ou entre quaisquer submódulos independentes do módulo principal;
            \item Módulos de software/firmware que são modificáveis;
            \item Módulos de software/firmware que não são modificáveis;
            \item Outros tipos de componentes que não estão listados acima.
        \end{itemize}
        \item \textbf{EN.III.1.1.02:} Verificar se a documentação especifica a fronteira criptográfica. A fronteira criptográfica deve incluir qualquer hardware ou software que insere, recebe, processa ou emite parâmetros de segurança importantes que poderiam conduzir ao comprometimento de informações sensíveis se não controlados adequadamente.
        \item \textbf{EN.III.1.1.03:} Verificar se todos os componentes de hardware, software e firmware dentro da fronteira criptográfica estão incluídos na “lista de componentes principais”, e se há componentes fora da fronteira criptográfica que não estão listados como componentes do módulo criptográfico.
        \item \textbf{EN.III.1.1.04:} Verificar se a documentação mostra explicitamente e precisamente onde o perímetro físico da fronteira criptográfica está situado, incluindo detalhes sobre os seus componentes. Além disso, analisar se a documentação contém uma lista das portas conectadas aos equipamentos externos, todos os fluxos de informação significativa e processamentos a serem desempenhados dentro da fronteira criptográfica, além de toda informação recebida e emitida.
        \item \textbf{EN.III.1.1.05:} Verificar se a fronteira criptográfica é fisicamente contínua, de tal forma que não haja lacunas que possam permitir entrada, saída ou outro tipo de acesso não controlado ao módulo criptográfico. O projeto do módulo criptográfico deve também assegurar que não tenham interfaces não controladas para o interior ou para fora do módulo criptográfico que possam passar PCSs (Parâmetros Críticos de Segurança), dados em texto legível, ou outras informações que, se mal utilizadas ou utilizadas de forma inadequada, possam conduzir a um comprometimento da segurança.
        \item \textbf{EN.III.1.1.06:} Verificar se todos os componentes que estão identificados no diagrama de blocos pertencem à fronteira criptográfica.
        \item \textbf{EN.III.1.1.07:} Verificar se a documentação descreve os componentes de software/firmware executados pelo módulo criptográfico, bem como seus serviços desempenhados e os dispositivos de memória que armazenam dados e o código executável.
        \item \textbf{EN.III.1.1.08:}  Analisar a documentação e identificar os componentes de hardware internos ou externos ao módulo criptográfico que interagem com o processador listado na “lista de componentes principais”
    \end{itemize}

    \vspace{1cm} 
    
    \subsubsection*{EN.III.1.1.01 - Parecer:}
    
    Mauris at pretium dolor. Sed ac magna eu lacus congue placerat. Vivamus volutpat, mi nec venenatis pharetra, massa massa eleifend elit, vel malesuada purus tortor sed velit. Phasellus fringilla facilisis semper. Nunc accumsan aliquet malesuada. Sed at hendrerit enim. Suspendisse sed leo at orci tristique eleifend. Mauris placerat, quam ut vestibulum fermentum, nisl est molestie nunc, et tempus quam eros vitae sapien  conforme figura \ref{fig:minha-evidencia-1}
    
    \vspace{1cm} 
    
    \subsubsection*{EN.III.1.1.02 - Parecer:}
    
    Mauris at pretium dolor. Sed ac magna eu lacus congue placerat. Vivamus volutpat, mi nec venenatis pharetra, massa massa eleifend elit, vel malesuada purus tortor sed velit. Phasellus fringilla facilisis semper. Nunc accumsan aliquet malesuada. Sed at hendrerit enim. Suspendisse sed leo at orci tristique eleifend. Mauris placerat, quam ut vestibulum fermentum, nisl est molestie nunc, et tempus quam eros vitae sapien  conforme figura \ref{fig:minha-evidencia-1}
    
    \vspace{1cm} 
    
    \subsubsection*{EN.III.1.1.03 - Parecer:}
    
    Mauris at pretium dolor. Sed ac magna eu lacus congue placerat. Vivamus volutpat, mi nec venenatis pharetra, massa massa eleifend elit, vel malesuada purus tortor sed velit. Phasellus fringilla facilisis semper. Nunc accumsan aliquet malesuada. Sed at hendrerit enim. Suspendisse sed leo at orci tristique eleifend. Mauris placerat, quam ut vestibulum fermentum, nisl est molestie nunc, et tempus quam eros vitae sapien  conforme figura \ref{fig:minha-evidencia-1}
    
    \vspace{1cm} 
    
    \subsubsection*{EN.III.1.1.04 - Parecer:}
    
    Mauris at pretium dolor. Sed ac magna eu lacus congue placerat. Vivamus volutpat, mi nec venenatis pharetra, massa massa eleifend elit, vel malesuada purus tortor sed velit. Phasellus fringilla facilisis semper. Nunc accumsan aliquet malesuada. Sed at hendrerit enim. Suspendisse sed leo at orci tristique eleifend. Mauris placerat, quam ut vestibulum fermentum, nisl est molestie nunc, et tempus quam eros vitae sapien  conforme figura \ref{fig:minha-evidencia-1}
    \vspace{1cm} 
    
    \subsubsection*{EN.III.1.1.05 - Parecer:}
    
    Mauris at pretium dolor. Sed ac magna eu lacus congue placerat. Vivamus volutpat, mi nec venenatis pharetra, massa massa eleifend elit, vel malesuada purus tortor sed velit. Phasellus fringilla facilisis semper. Nunc accumsan aliquet malesuada. Sed at hendrerit enim. Suspendisse sed leo at orci tristique eleifend. Mauris placerat, quam ut vestibulum fermentum, nisl est molestie nunc, et tempus quam eros vitae sapien  conforme figura \ref{fig:minha-evidencia-1}
    \vspace{1cm} 
    
    \subsubsection*{EN.III.1.1.06 - Parecer:}
    
    Mauris at pretium dolor. Sed ac magna eu lacus congue placerat. Vivamus volutpat, mi nec venenatis pharetra, massa massa eleifend elit, vel malesuada purus tortor sed velit. Phasellus fringilla facilisis semper. Nunc accumsan aliquet malesuada. Sed at hendrerit enim. Suspendisse sed leo at orci tristique eleifend. Mauris placerat, quam ut vestibulum fermentum, nisl est molestie nunc, et tempus quam eros vitae sapien  conforme figura \ref{fig:minha-evidencia-1}
    \vspace{1cm} 
    
    \subsubsection*{EN.III.1.1.07 - Parecer:}
    
    Mauris at pretium dolor. Sed ac magna eu lacus congue placerat. Vivamus volutpat, mi nec venenatis pharetra, massa massa eleifend elit, vel malesuada purus tortor sed velit. Phasellus fringilla facilisis semper. Nunc accumsan aliquet malesuada. Sed at hendrerit enim. Suspendisse sed leo at orci tristique eleifend. Mauris placerat, quam ut vestibulum fermentum, nisl est molestie nunc, et tempus quam eros vitae sapien  conforme figura \ref{fig:minha-evidencia-1}
    
    \vspace{1cm} 
    
    \subsubsection*{EN.III.1.1.08 - Parecer:}
    
    Mauris at pretium dolor. Sed ac magna eu lacus congue placerat. Vivamus volutpat, mi nec venenatis pharetra, massa massa eleifend elit, vel malesuada purus tortor sed velit. Phasellus fringilla facilisis semper. Nunc accumsan aliquet malesuada. Sed at hendrerit enim. Suspendisse sed leo at orci tristique eleifend. Mauris placerat, quam ut vestibulum fermentum, nisl est molestie nunc, et tempus quam eros vitae sapien  conforme figura \ref{fig:minha-evidencia-1}
    
    \vspace{1cm} 
    
    
    % %%%%%%%%%%%%%%%%%%%%%%%%%%%%%%%%%%%%%%%%%%%%%%%%%%%%%%%
    % --- CÓDIGO CORRETO PARA INSERIR FIGURAS ---
    % %%%%%%%%%%%%%%%%%%%%%%%%%%%%%%%%%%%%%%%%%%%%%%%%%%%%%%%
    \begin{figure}[htbp] % <-- O [htbp] é a mágica!
        \centering % Centraliza a figura
        
        % --- !!! MUDE O NOME E A LARGURA !!! ---
        % Use width=\textwidth para largura máxima ou 0.8\textwidth, etc.
        \includegraphics[width=0.9\textwidth]{figures/circuit.png}
        
        % --- !!! ADICIONE UMA LEGENDA !! ---
        % Isto é ESSENCIAL para a Lista de Figuras funcionar
        \caption{Esta é a legenda da sua figura (ex: Evidência do Ensaio EN.III.1.1.01).}
        
        % Opcional, mas boa prática para referenciar no texto
        \label{fig:minha-evidencia-1} 
    \end{figure}
    % %%%%%%%%%%%%%%%%%%%%%%%%%%%%%%%%%%%%%%%%%%%%%%%%%%%%%%%
    % --- FIM DO CÓDIGO DA FIGURA ---
    
    % %%%%%%%%%%%%%%%%%%%%%%%%%%%%%%%%%%%%%%%%%%%%%%%%%%%%%%%
    % --- CÓDIGO CORRETO PARA INSERIR TABELAS ---
    % %%%%%%%%%%%%%%%%%%%%%%%%%%%%%%%%%%%%%%%%%%%%%%%%%%%%%%%
    
    % 'h!' tenta forçar a tabela a ficar "aqui"
    \begin{table}[h!]
        \centering
        \begin{tabular}{ll}
            \toprule
            \textbf{Unidade (X)} & \textbf{Desvio padrão} \\
            \midrule
             XX & XX\\
             XX & XX \\  
             XX & XX \\  
             XX & XX \\  
             XX & XX \\  
            \bottomrule
        \end{tabular}
        \caption{Tabela de exemplo}
        \label{table:Tabela3}
    \end{table}
    
    % %%%%%%%%%%%%%%%%%%%%%%%%%%%%%%%%%%%%%%%%%%%%%%%%%%%%%%%
    % --- FIM DO CÓDIGO PARA INSERIR TABELAS ---
    % %%%%%%%%%%%%%%%%%%%%%%%%%%%%%%%%%%%%%%%%%%%%%%%%%%%%%%%
    
    % %%%%%%%%%%%%%%%%%%%%%%%%%%%%%%%%%%%%%%%%%%%%%%%%%%%%%%%
    % --- CÓDIGO CORRETO PARA INSERIR SNIPPET DE CÓDIGO-FONTE
    % %%%%%%%%%%%%%%%%%%%%%%%%%%%%%%%%%%%%%%%%%%%%%%%%%%%%%%%
    
    
    \begin{lstlisting}[language=Python, caption={Exemplo de snippet de código}]
    
    import requests
    
    def check_api_status(url):
        try:
            r = requests.get(url + "/api/status")
            if r.status_code == 200:
                print("API Status: OK")
                print(r.json())
            else:
                print("API Status: ERROR")
        except Exception as e:
            print("Failed to connect: " + str(e))
    
    check_api_status("http://192.168.0.1")
    \end{lstlisting}
    
    % %%%%%%%%%%%%%%%%%%%%%%%%%%%%%%%%%%%%%%%%%%%%%%%%%%%%%%%
    % ---fIM DO CÓDIGO CORRETO PARA INSERIR SNIPPET 
    % %%%%%%%%%%%%%%%%%%%%%%%%%%%%%%%%%%%%%%%%%%%%%%%%%%%%%%%
    
    
    
    \vspace{0.5cm}
    
    
    
    \vspace{2cm} % Espaço para o parecer
    
    \newpage


\subsection{REQUISITO III.1.2} 
    
    % --- Tabela minimalista (Design Limpo) ---
    
    \begin{tabularx}{\textwidth}{l X} 
    \toprule
    
    \textbf{Referência Normativa} & \textbf{Requisito e Item} \\ 
    \midrule
    
    \rowcolor{gray!15} 
    
    \textbf{MCT 7 Vol. II} & 
    "A parte interessada deve fornecer documentação específica que descreve a configuração física do módulo criptográfico." \\
    \bottomrule
    \end{tabularx}
    % --- Fim da Tabela ---
    
    \vspace{1cm} 
    
    % --- Procedimentos (Texto Livre) ---
    \subsubsection*{Procedimentos de ensaio para NSH 1, 2 e 3:}
    \begin{itemize}[leftmargin=*, nosep]
        \item \textbf{EN.III.1.2.01:} Analisar a documentação e verificar se a identificação do módulo criptográfico está enquadrada numa das seguintes definições, conforme padrão FIPS PUB 140-2, Seção 4.5: single-chip module, multi-chip embedded module ou multi-chip standalone module.
        \item \textbf{EN.III.1.2.02:} Analisar a documentação e verificar a disposição interna (internal layout) do módulo criptográfico por meio de desenhos técnicos, esboços ou diagramas de blocos que identifiquem cada bloco dos componentes de hardware.
        \item \textbf{EN.III.1.2.03:} Analisar a documentação e verificar as principais montagens/encapsulamentos físicos do
    módulo e como são dispostas tais montagens/encapsulamentos no módulo criptográfico.
        \item \textbf{EN.III.1.2.04:}Analisar a documentação e verificar a descrição dos parâmetros físicos principais do módulo
    criptográfico, constando, no mínimo, dos seguintes itens:
        \begin{itemize}
            \item Forma de encapsulamento/revestimento e dimensões aproximadas, incluindo quaisquer
    interfaces de acesso ou coberturas/camadas;
            \item Dimensões, disposição (layout) e interconexões de placa(s) de circuito impresso;
            \item Localização da fonte de alimentação de energia, conversores de energia e entradas e
    saídas de energia;
            \item Ativação de componentes interconectados por meio de condutores elétricos
    (interconection wiring runs): rotas e terminais;
            \item Arranjos de refrigeração, tais como, pratos de condução (conduction plates), duto de
    refrigeração (cooling airflow), trocadores de calor (heat exchanger), palhetas de
    refrigeração (cooling fins), ventiladores (fans), ou outros arranjos para a remoção de calor
    do módulo;
            \item Outros tipos de componentes não listados acima.
        \end{itemize}
        \item \textbf{EN.III.1.2.05:}Verificar se os itens descritos no EN.III.1.2.04 estão em conformidade com as estruturas físicas contidas na "lista de componentes principais".
    \end{itemize}  
    
    \vspace{1cm} 
    
    \subsubsection*{EN.III.1.2.01 - Parecer:}
    
    \vspace{1cm} 
    
    \subsubsection*{EN.III.1.2.02 - Parecer:}
    
    \vspace{1cm} 
    
    \subsubsection*{EN.III.1.2.03 - Parecer:}
   
    \vspace{1cm} 
    
    \subsubsection*{EN.III.1.2.04 - Parecer:}
    
    
    \vspace{1cm} 
    
    \subsubsection*{EN.III.1.2.05 - Parecer:}
    
    
    
    \newpage

\end{document}